\subsection{Datasets and preprocessing}
\owner{Member 3; specification aligned with the supplied Member 4 code.}
We use Fashion-MNIST \citep{xiao2017fashion}, CIFAR-10 \citep{krizhevsky2009}, and
Tiny Shakespeare \citep{karpathychar}. For each image task, a permutation with split seed
6704 selects training examples, followed by a disjoint validation subset from the official
training data. A separate permutation from the same generator selects the official-test
subset. These are random subsets, not stratified samples. Full indices and SHA-256 split
fingerprints are saved in \texttt{data/manifests/}. Fashion inputs use mean/std
$0.2860/0.3530$; CIFAR channels use means $(0.4914,0.4822,0.4465)$ and standard deviations
$(0.2470,0.2435,0.2616)$. There is no augmentation.

Text boundaries are character offsets 892,312 and 1,003,851. Each contiguous split is
windowed separately into 64-character inputs and one-character-shifted targets, so windows
cannot cross boundaries. The sorted vocabulary is obtained from the entire source text,
as in Member 4's code; this uses held-out character inventory but no held-out loss for fitting.
This choice and overlapping evaluation windows limit the interpretation of generalization.

\begin{table}[H]\centering\small
\caption{Actual common task specification (Member 3). Text sizes count possible windows.}
\label{tab:data}
\begin{tabularx}{\linewidth}{lXrr}\toprule
Task (writer) & Model & Train / validation / test & Parameters\\\midrule
Fashion (M3) & ReLU MLP $784\to128\to64\to10$ & 6,000 / 1,000 / 1,000 & 109,386\\
CIFAR (M3) & Two-convolution CNN & 5,000 / 1,000 / 1,000 & 25,578\\
Text (M3) & Character GRU & 892,248 / 111,475 / 111,479 & 25,121\\
\bottomrule\end{tabularx}\end{table}

\subsection{Model definitions and executed protocol}
\owner{Member 3; exact counts are checked against the generated dataset summary.}
The CNN uses $3\times3$ convolutions with padding 1, channels $3\to16\to32$,
ReLU and $2\times2$ max pooling after each convolution, and a $32\cdot8\cdot8\to10$ classifier.
The GRU has a 65-character vocabulary, embedding dimension 32, one hidden layer of width 64,
and a 64-to-65 output projection. All parameters are unconstrained; loss is cross-entropy
(CE), averaged over images or target characters. No model uses dropout or batch normalization.

Main training uses 300 gradient calls, batches of 64 images or 32 text windows, and seeds
1, 2, 3. Batches are sampled with replacement using seed $100000+\text{training seed}$.
Within each task/seed pair, initialization and sampled batch indices match exactly, as
verified by hashes. ADOPT-0 makes 299 parameter updates because its first call initializes
$v$. Validation-only tuning uses seed 11, 100 calls per trial and
$\alpha\in\{0.0003,0.001,0.003\}$ for each optimizer/task. The minimum final validation CE
selects the learning rate before any main test evaluation.

Evaluation occurs at calls 0, 50, \ldots, 300, using the first 1,024 examples of each split
(all 1,000 examples for image validation/test). For text these are overlapping windows,
covering only a short prefix of each held-out region. Training CE also uses this fixed prefix,
not all training examples. Each final model is tested once. Times include training and scheduled
train/validation evaluation, but exclude downloads, preprocessing and final test evaluation;
the logs additionally retain total time including the test. All runs use an Apple M1 Pro,
16 GiB RAM, CPU only, four PyTorch threads, Python 3.12.14 and PyTorch 2.14.1.
These are aligned settings, not proof of identical Member 4 runs: that member's actual
manifests, framework version and evaluation-call overrides were not supplied.
